\section{Official Artifacts}

\subsection{Primary Sources Read for the Thesis}
The final thesis was written from the implemented repository and official artifacts:
\begin{itemize}
  \item \texttt{README.md}
  \item \texttt{docs/ARCHITECTURE.md}
  \item \texttt{docs/REPRODUCIBILITY.md}
  \item the preserved current thesis source under \texttt{thesis/original/}
  \item \texttt{artifacts/index/current.json}
  \item \texttt{artifacts/evaluation/golden\_benchmark\_100\_extended.jsonl}
  \item \texttt{artifacts/evaluation/final/ablation\_retrieval\_final\_report.md}
  \item \texttt{artifacts/evaluation/final/ablation\_retrieval\_final\_results.json}
  \item \texttt{artifacts/evaluation/final/end\_to\_end\_final\_report.md}
  \item \texttt{artifacts/evaluation/final/end\_to\_end\_final\_results.json}
\end{itemize}

\subsection{Artifact Availability Note}
The current checkout contains the official index manifest and final evaluation outputs. The reproducibility guide references generated chunk and graph paths such as \path{artifacts/chunks/} and \path{artifacts/graph/}; those directories are not present in this checkout. The final thesis therefore uses \path{artifacts/index/current.json} for corpus/index counts and \path{artifacts/evaluation/final/} for evaluation metrics.

\section{Metric Formulas}

\subsection{Retrieval Metrics}
For an in-corpus query $q$, let $G(q)$ be the required citation set and let $\mathcal{C}_{10}(q)$ be the citation coordinates retrieved in the top-10 context window. Recall@10 is:
\[
\operatorname{Recall}@10(q) =
\frac{|G(q) \cap \mathcal{C}_{10}(q)|}{|G(q)|}.
\]
The reported Recall@10 is the macro-average over 94 in-corpus queries. Required Citation Coverage is reported with the same citation-level coverage definition in the official ablation report.

For forbidden citations, let $F(q)$ be the benchmark-defined forbidden set. The violation rate is:
\[
\operatorname{ForbiddenViolationRate} =
\frac{1}{94}
\sum_{q}
\mathbb{I}\left[F(q) \cap \mathcal{C}_{10}(q) \neq \emptyset\right].
\]

\subsection{Adjusted End-to-End Score}
The 100-query benchmark contains 94 in-corpus queries and 6 out-of-corpus refusal tests. For in-corpus queries, a pass requires retrieval, answer, citation, and quality checks. For out-of-corpus queries, a pass requires refusal or insufficient-context behavior. The adjusted end-to-end score is:
\[
\operatorname{AdjustedE2E} =
\frac{
\sum_{q \in Q_{\mathrm{in}}}\mathbb{I}[\operatorname{Pass}_{\mathrm{in}}(q)]
+
\sum_{q \in Q_{\mathrm{ooc}}}\mathbb{I}[\operatorname{RefusalPass}(q)]
}{100}.
\]
The final official score is $0.780$.

\section{Reproducibility Commands}

\subsection{Retrieval Ablation}
\begin{ThesisCode}
.venv\Scripts\python.exe scripts\ablation_retrieval_100.py `
  --benchmark-path artifacts\evaluation\golden_benchmark_100_extended.jsonl `
  --output-prefix ablation_retrieval_100_final_candidate
\end{ThesisCode}

\subsection{End-to-End Evaluation}
\begin{ThesisCode}
.venv\Scripts\python.exe scripts\evaluate_end_to_end_rag.py `
  --benchmark-path artifacts\evaluation\golden_benchmark_100_extended.jsonl `
  --output-prefix end_to_end_100_final_candidate
\end{ThesisCode}

\section{Representative Benchmark Cases}

\subsection{Successful In-Corpus Case}
Query: \textit{Female retirement age in 2026.}

Required evidence: Labor Code retirement-age rule and Decree 135/2020/ND-CP table evidence.

Outcome: pass. The system retrieves the relevant retirement-age evidence and returns a citation-grounded answer.

\subsection{Successful Out-of-Corpus Case}
Query: \textit{Minimum wage by region in 2026.}

Expected behavior: refuse or report insufficient indexed context.

Outcome: pass. The required legal instrument is outside the indexed corpus, so the system does not provide an unsupported answer.

\subsection{Remaining Failure Case}
Query: \textit{The company says there is little work and transfers me to another job for nearly three months without asking me; is that acceptable?}

Required evidence: Labor Code Article 29 clauses on temporary transfer.

Outcome: fail. The final report marks the case as a missing-required-context retrieval failure.
